% -----------------------------------------------------------
%  Anexo --- Casos de Uso restantes
% -----------------------------------------------------------

\renewcommand{\cutop}[2]{%
  \multicolumn{2}{|l|}{\textbf{#1\;\;#2}}\\\hline
}
\renewcommand{\cusec}[2]{%
  \multicolumn{2}{|p{0.96\linewidth}|}{\textbf{#1:}\;#2}\\\hline
}
\renewcommand{\cuheader}[2]{\textbf{#1} & \textbf{#2}\\\hline}

Este anexo consolida los casos de uso que quedaron fuera del alcance operativo de la version~1.0 de \textit{DOMU}. Se mantienen como funcionalidades proyectadas para una version posterior del sistema y se preservan aqui para no perder la trazabilidad funcional del informe.

\begin{longtable}{|p{0.09\textwidth}|p{0.24\textwidth}|p{0.11\textwidth}|p{0.12\textwidth}|p{0.34\textwidth}|}
\caption{Casos de uso restantes planificados para versiones futuras.}\label{tab:casos-restantes}\\
\hline
\textbf{CU} & \textbf{Nombre} & \textbf{Actor principal} & \textbf{RF asociado} & \textbf{Alcance pendiente} \\
\hline
\endfirsthead
\hline
\textbf{CU} & \textbf{Nombre} & \textbf{Actor principal} & \textbf{RF asociado} & \textbf{Alcance pendiente} \\
\hline
\endhead
\hline
\endfoot

CU\_06 & Gestionar morosidad y conciliacion & Administrador & RF\_08 & Requiere conciliacion bancaria, regularizacion automatica de deuda, aplicacion formal de intereses y seguimiento administrativo integral de morosidad. \\
\hline
CU\_11 & Mantenimiento preventivo & Administrador & RF\_14 & Requiere planificacion de mantenciones, alertas de vencimiento, evidencia de ejecucion y reprogramacion de actividades preventivas. \\
\hline
CU\_12 & Gestionar estacionamientos (titulares y pases) & Administrador & RF\_09 (version previa) & Requiere administracion de cupos, pases temporales, disponibilidad y reglas de uso para residentes y visitas. \\
\hline
CU\_18 & Registrar mantencion de ascensor en bitacora & Conserje & RF\_14 & Requiere una bitacora especializada para ascensores, revisiones tecnicas y relacion con planes de mantencion futura. \\
\hline
\end{longtable}

% CU_06
\begin{table}[H]\centering
\begin{tabular}{|p{0.475\linewidth}|p{0.475\linewidth}|}\hline
\cutop{CU\_06}{Gestionar morosidad y conciliacion}
\cusec{Precondiciones}{Administrador autenticado; estados de cuenta y comprobantes disponibles para revision financiera.}
\cusec{Includes}{---}
\cusec{Extends}{---}
\cuheader{Administrador}{Sistema}
1) Accede al modulo de morosidad. & 2) Muestra deudores, saldos pendientes y movimientos asociados. \\\hline
3) Revisa pagos pendientes y aplica acciones administrativas. & 4) Calcula intereses, registra observaciones y prepara la conciliacion. \\\hline
5) Importa o cruza informacion bancaria. & 6) Sugiere coincidencias, detecta diferencias y actualiza la regularizacion cuando corresponda. \\\hline
\multicolumn{2}{|l|}{\textbf{Flujo de eventos alternativos}}\\\hline
5a. El archivo bancario es invalido o no coincide con el periodo. & 5b. Rechaza la importacion y solicita una nueva fuente de datos. \\\hline
\cusec{Post condiciones}{La morosidad queda conciliada o actualizada con trazabilidad financiera.}
\end{tabular}
\caption{CU\_06 --- Gestionar morosidad y conciliacion}
\label{tab:cu06}
\end{table}

% CU_11
\begin{table}[H]\centering
\begin{tabular}{|p{0.475\linewidth}|p{0.475\linewidth}|}\hline
\cutop{CU\_11}{Mantenimiento preventivo}
\cusec{Precondiciones}{Administrador autenticado; activos o equipos registrados en la comunidad.}
\cusec{Includes}{---}
\cusec{Extends}{---}
\cuheader{Administrador}{Sistema}
1) Define un plan preventivo para un activo o instalacion. & 2) Registra periodicidad, responsable y fecha objetivo. \\\hline
3) Programa la actividad de mantencion. & 4) Genera alertas de vencimiento y seguimiento del plan. \\\hline
5) Registra evidencia de ejecucion o reprogramacion. & 6) Actualiza el historial preventivo del activo. \\\hline
\multicolumn{2}{|l|}{\textbf{Flujo de eventos alternativos}}\\\hline
5a. La mantencion no puede ejecutarse en la fecha programada. & 5b. Reagenda la actividad y registra el motivo. \\\hline
\cusec{Post condiciones}{El plan preventivo queda programado, ejecutado o reprogramado con evidencia.}
\end{tabular}
\caption{CU\_11 --- Mantenimiento preventivo}
\label{tab:cu11}
\end{table}

% CU_12
\begin{table}[H]\centering
\begin{tabular}{|p{0.475\linewidth}|p{0.475\linewidth}|}\hline
\cutop{CU\_12}{Gestionar estacionamientos (titulares y pases)}
\cusec{Precondiciones}{Administrador autenticado; mapa y reglas de estacionamientos configuradas.}
\cusec{Includes}{---}
\cusec{Extends}{---}
\cuheader{Administrador}{Sistema}
1) Accede al modulo de estacionamientos. & 2) Muestra cupos, titulares y disponibilidad actual. \\\hline
3) Asigna o modifica un titular. & 4) Valida disponibilidad, vigencia y reglas de uso. \\\hline
5) Emite un pase temporal o reasigna un cupo. & 6) Registra la autorizacion y controla su vigencia. \\\hline
\multicolumn{2}{|l|}{\textbf{Flujo de eventos alternativos}}\\\hline
4a. El cupo seleccionado ya se encuentra ocupado. & 4b. Informa el conflicto y solicita elegir otra alternativa. \\\hline
\cusec{Post condiciones}{La asignacion o pase queda registrado y trazable para control futuro.}
\end{tabular}
\caption{CU\_12 --- Gestionar estacionamientos (titulares y pases)}
\label{tab:cu12}
\end{table}

% CU_18
\begin{table}[H]\centering
\begin{tabular}{|p{0.475\linewidth}|p{0.475\linewidth}|}\hline
\cutop{CU\_18}{Registrar mantencion de ascensor en bitacora}
\cusec{Precondiciones}{Conserje autenticado; existe una revision tecnica o mantencion del ascensor en curso o finalizada.}
\cusec{Includes}{---}
\cusec{Extends}{CU\_11 Mantenimiento preventivo, cuando la revision forma parte de un plan formal.}
\cuheader{Conserje}{Sistema}
1) Accede a la bitacora tecnica del ascensor. & 2) Muestra formulario de registro de evento o mantencion. \\\hline
3) Ingresa responsable, hallazgos, fecha y observaciones. & 4) Registra la entrada y la vincula al equipo o servicio correspondiente. \\\hline
5) Adjunta evidencia si existe. & 6) Guarda la bitacora y notifica a administracion para seguimiento. \\\hline
\multicolumn{2}{|l|}{\textbf{Flujo de eventos alternativos}}\\\hline
5a. No existe evidencia adjunta. & 5b. Permite guardar el registro con solo antecedentes descriptivos. \\\hline
\cusec{Post condiciones}{La mantencion queda registrada en una bitacora tecnica especializada del ascensor.}
\end{tabular}
\caption{CU\_18 --- Registrar mantencion de ascensor en bitacora}
\label{tab:cu18}
\end{table}

La postergacion de estos casos de uso permitio enfocar la version~1.0 en los modulos operativos ya entregados: control de accesos, encomiendas, incidentes, reservas, finanzas, tareas, comunicaciones comunitarias, biblioteca y gobernanza basica.
